\newpage
\section{Group Cohomology and Subgroup}
When $L/K$ is a finite Galois extension with group $G$, we have $H^1(G, L^\times) = 0$.
What about $H^1(G, L)$?
Before we answer this question, we take a look at general theory of cohomology groups.

\subsection{Induced and Coinduced Modules}
Let's consider more general setting.
Let $\alpha: R \to S$ be a ring homomorphism, and they don't have to be commutative.
Suppose that $S$ is a projective(or free) $R$-module via $\alpha$.
Let $M$ be an $R$-modules, and let $N, P$ be $S$-modules($R$-modules as well), and let $h: M \to N$ be an $R$-linear map.

\begin{definition}[Induced and Coinduced Modules]
    \noindent\begin{enumerate}
        \item The \textit{induced $S$-module} $\Ind_R^S(M)$ is defined as $S \otimes_R M$.
        \item The \textit{coinduced $S$-module} $\Coind_R^S(M)$ is defined as $\Hom_R(S, M)$.
        We note that
        for $f \in \Hom_R(S, N)$ and $s \in S$, we define $(s \cdot f)(t) = f(st)$ for $t \in S$.
    \end{enumerate}
\end{definition}

\begin{proposition}
    There is a canonical isomorphism $$\Ext_S^i(N, \Hom_R(S, M)) \cong \Ext_R^i(N, M).$$
\end{proposition}
\begin{proof}
    Let $P^\bullet \to N$ be a projective(or free resp.) resolution of $P$ as an $S$-module.
    Then it is also a projective resolution of $N$ as an $R$-module since $S$ is a projective $R$-module.
    So we have
    \begin{align*}
        \Ext_R^i(N, M) &\cong H^i(\Hom_R(P^\bullet, M)) \\
        &\cong H^i(\Hom_S(P^\bullet, \Hom_R(S, M))) \\
        &\cong \Ext_S^i(N, \Hom_R(S, M)).
    \end{align*}
\end{proof}

\subsection{Group algebra and subgroup algebra}

\begin{lemma}
    Suppose that $G$ is a discrete group and $H\hookrightarrow G$ is a subgroup of finite index.
    Then we have that $\Z[G]$ is a free $\Z[H]$-module with basis $G/H$.
\end{lemma}
\begin{proof}
    We know that $\Z[H]$ is a free $\Z$-module with basis $H$,
    and $\Z[G]$ is a free $\Z$-module with basis $G$.
    Let $g_0 = e, g_1, \ldots, g_{n-1}$ be a set of representatives where $n = [G:H]$.

    Consider a map $\Phi: \bigoplus_{i=0}^{n-1} \Z[H]\tilde{g_i} \to \Z[G]$ defined by $(h_0, h_1, \ldots, h_{n-1}) \mapsto \sum_{i=0}^{n-1} g_i h_i$.
    It is clear $\Phi$ is surjective since every $g = hg_i$ for some $h \in H$.
    We claim that $\Phi$ is injective as well.
    With $g \mapsto h \tilde{g_i}$, we get a map $\Psi: \Z[G] \to \bigoplus_{i=0}^{n-1} \Z[H]\tilde{g_i}$.
    We notice that $\Phi \circ \Psi = \mathrm{id}_{\Z[G]}$, so $\Phi$ is injective.
\end{proof}

So we can consider $R = \Z[H]$ and $S = \Z[G]$ in the previous subsection.
We can consider the induced and coinduced module maps: $\Ind_H^G:=\Ind_R^S$ and $\Coind_H^G:=\Coind_R^S$.
We can set $M$ to be an $H$-module, and $N$ to be $\Z$ as a trivial $G$-module.
So the following proposition is a direct consequence of the previous proposition.

\begin{proposition}[Shapiro's Lemma]
    $$\Ext_S^i(\Z, \Coind^G_H(A)) \cong \Ext_R^i(\Z, A).$$
\end{proposition}

\begin{corollary}
    Let $G$ be a finite group and $A$ be any abelian group.
    Then we have $H^i(G, \Coind^G_*(A)) = A$ if $i = 0$ and $0$ otherwise
    where $\Coind^G_*(A) = \Hom_\Z(\Z[G], A)$.
\end{corollary}
\begin{proof}
    It follows from Shapiro's lemma when $H = \{e\}$.
\end{proof}

\subsection{More Discussion on Hilbert 90}
\begin{theorem}
    Let $L/K$ be a finite Galois extension with group $G$.
    Then we have $H^0(G, L) = K$ and $H^i(G, L) = 0$ for all $i > 0$.
\end{theorem}
\begin{proof}
    By normal basis theorem, we have some $x \in L$ such that $\{g(x)\}_{g \in G}$ is a basis of $L$ over $K$.
    We observe $\Hom_\Z(\Z[G], K) \approx \Z[G]\underset{\Z}{\otimes} K$
    and $K[G] \cong L$ since $\Z[G]$ is a free $\Z$-module.
    Hence $\Coind^G_*(K) \cong L$.
    From the previous corollary, we have $H^i(G, L) = H^i(G, \Coind^G_*(K)) = K$ if $i = 0$ and $0$ otherwise.
\end{proof}

\begin{corollary}
    If $L/K$ is a cyclic extension with group $G = \langle \sigma \rangle$, then we have
    an exact sequence \(0 \to K \to L \xrightarrow{\sigma - 1} L \xrightarrow{\Tr_{L/K}} K \to 0\).
\end{corollary}
\begin{proof}
    The exactness at the first $L$ is clear since $\sigma(x) = x$ if and only if $x \in K$.
    We have $0=H^1(G,L) = {}_NL /(\sigma - 1)L$,
    and the norm map is exactly the trace map $\Tr_{L/K}$;
    hence we have the exactness at the second $L$.
    We also have $0 = H^2(G, L) = L^G / NL = K/\Tr_{L/K}(L)$; we have the exactness at the last $K$.
\end{proof}
